
\makeabstract
{ % #1: Title
Using GCamP6s Fluorescence in Zebrafish Hair Cells to Understand the Effects of Glutamate Agonist AMPA on Hair Cell Activity 
}
{ % #2: Authorship
Sana Farhat\textsuperscript{1}, 
Stacey Beganny\textsuperscript{1}, 
Josef G. Trapani\textsuperscript{1}, 
}
{ % #3: Affiliations
\textsuperscript{1} Neuroscience Program, Amherst College, Amherst, MA \\
}
{ % #4: Purpose
Glutamate excitotoxicity refers to the neuronal damage and cell death associated with the accumulation of excess glutamate. Studies suggest that the accumulation of excess glutamate damages sensory hair cells, which are essential for hearing and balance. This study aimed to further assess the effects of alpha-amino-3-hydroxy-5-methyl-4-isoxazolepropionic acid (AMPA), a glutamate agonist, on hair cell activity in zebrafish. Because zebrafish hair cells activate differently to anterior-to-posterior (A \textrightarrow{} P) and posterior-to-anterior ( P \textrightarrow{} A) fluid-jet stimulation, this study investigated the effects of AMPA on hair cell activity during A \textrightarrow{} P and P \textrightarrow{} A stimulation. 
}
{ % #5: Methods
Using calcium imaging procedures with zebrafish, this study (1) validated hair cell responsiveness to A \textrightarrow{} P and P \textrightarrow{} A stimulation, (2) confirmed that hair cells respond differentially to A \textrightarrow{} P and P \textrightarrow{} A stimulation and classified hair cells as either A \textrightarrow{} P or P \textrightarrow{} A responding, and (3) collected baseline measures and monitored the effects of AMPA in 0.1\% dimethyl sulfoxide (DMSO) in E3 over 50 minutes. .
}
{ % #6: Results
The results suggest that AMPA in 0.1\% DMSO in E3 appears to inhibit hair cell activity during A \textrightarrow{} P stimulation for both A \textrightarrow{} P and P \textrightarrow{} A responding hair cells, but not during P \textrightarrow{} A stimulation for either group. Spontaneously, this study also discovered that 0.4\% DMSO in E3 appears to completely inhibit general hair cell activity during A \textrightarrow{} P stimulation. 
}
{ % #7: Conclusion
These results may suggest that AMPA affects hair cell activity in a direction-specific way. However, because 0.4\% DMSO in E3 also inhibited hair cell activity, future studies must assess whether DMSO also inhibits hair cell activity at 0.1\%, and whether it contributes to the effects of AMPA in 0.1\% DMSO in E3. Studies should administer DMSO at 0.1\% in E3 alone as well as AMPA in E3 without DMSO. Furthermore, both DMSO and AMPA groups must be compared to a pure control group with E3 only. A limitation of this study is that the A \textrightarrow{} P stimulation was stronger than the P \textrightarrow{} A stimulation. Future studies must assess whether the direction-specific effects of AMPA remain even after using similar strength of A \textrightarrow{} P and P \textrightarrow{} A stimulation. 
}
 \newpage